\documentclass[../../main.tex]{subfiles}
\begin{document}

{\bf [解答]}

\begin{figure}[htb]
  \centering
  \begin{tikzpicture}[scale=1.05, >=stealth,
      circleline/.style={gray!70, thin},
      chord/.style={blue!75!black, thick},
    auxiliary/.style={teal!70!black, thick}]
    \coordinate (O) at (0,0);
    \coordinate (C) at (-2,0);
    \coordinate (D) at (2,0);
    % A, B are on the circle of radius 2. P is the exact chord intersection.
    \coordinate (A) at (0.5,-1.936492);
    \coordinate (B) at (1.3,1.519868);
    \path[name path=diameter] (C) -- (D);
    \path[name path=chordAB] (A) -- (B);
    \path[name intersections={of=diameter and chordAB, by=P}];

    \draw[circleline] (O) circle (2);
    \draw[circleline, thick] (C) -- (D);
    \draw[chord] (C) -- (A) -- (D);
    \draw[auxiliary] (C) -- (B) -- (D);
    \draw[chord] (A) -- (B);

    \node[circle, fill=gray!75, inner sep=1.2pt] at (O) {};
    \node[circle, fill=blue!75!black, inner sep=1.5pt] at (A) {};
    \node[circle, fill=teal!70!black, inner sep=1.5pt] at (B) {};
    \node[circle, fill=red!75!black, inner sep=1.6pt] at (P) {};
    \node[above] at (O) {$O$};
    \node[left] at (C) {$C$};
    \node[right] at (D) {$D$};
    \node[below] at (A) {$A$};
    \node[above] at (B) {$B$};
    \node[above left] at (P) {$P$};

    \pic[draw=gray!70, angle radius=5mm] {right angle=C--A--D};
    \pic[draw=gray!70, angle radius=5mm] {right angle=C--B--D};
    \pic[draw=blue!75!black, angle radius=5mm] {angle = A--C--D};
    \pic[draw=teal!70!black, angle radius=5mm] {angle = D--C--B};
    \pic[draw=red!75!black, angle radius=5mm] {angle = C--P--A};
    \pic[draw=gray!70, angle radius=5mm] {angle = B--D--A};
    \node[below right] at (-1.65,-0.13) {$\alpha$};
    \node[above right] at (-1.65,0.13) {$\beta$};
    \node[below left] at (P) {$\theta$};
    \draw[gray!70, thin] (1.83,-0.12) -- (2.08,-0.42);
    \node[below right] at (2.08,-0.42) {$\pi-\alpha-\beta$};
  \end{tikzpicture}
  \caption{直径 $CD$ と弦 $AB$ の交点 $P$}
  \label{fig:1969_2_chords}
\end{figure}

定円の中心を $O$ とする．
円周角の定理から $\angle CAD = \angle CBD = \angle R$ だから，
\begin{align}
  \tan\alpha \tan\beta
  &= \dfrac{|AD|}{|CA|} \cdot \dfrac{|BD|}{|CB|}
\end{align}
である．ここで分母分子に$\sin(\alpha + \beta)=\sin(\pi-\alpha-\beta)$をかけて
\begin{align}
  \tan\alpha \tan\beta
  = \dfrac{\dfrac{1}{2} |AD||BD|\sin(\pi - \alpha - \beta)}{\dfrac{1}{2} |AC||BC|\sin(\alpha + \beta)}
  = \dfrac{\triangle ADB}{\triangle ACB} \label{eq:1}
\end{align}
である．

次に$\angle CPA = \angle BPD = \theta$ とおけばこれらの三角形の面積は
\begin{align}
  \dfrac{\triangle ADB}{\triangle ACB}
  = \dfrac{\dfrac{1}{2} |AB||DP| \sin\theta}{\dfrac{1}{2} |AB||CP| \sin\theta}
  = \dfrac{|DP|}{|CP|} \label{eq:2}
\end{align}
である．

\cref{eq:1,eq:2}より
\begin{align}
  \tan\alpha \tan\beta = \dfrac{\triangle ADB}{\triangle ACB} = \dfrac{|DP|}{|CP|}
\end{align}
となって題意は示された．

\newpage

\end{document}
